% ============================================================
%  Relatorio de Projeto: Podcastia
%  Compilar com: xelatex Relatorio_Podcastia.tex  (2 passes)
%  Ou carregar directamente para o Overleaf (overleaf.com)
% ============================================================
\documentclass[12pt, a4paper]{report}

% --- Encoding & Língua ---
\usepackage[utf8]{inputenc}
\usepackage[T1]{fontenc}
\usepackage[portuguese]{babel}

% --- Tipografia ---
\usepackage{lmodern}
\usepackage{microtype}

% --- Geometria da Página ---
\usepackage[
  top=2.5cm, bottom=2.5cm,
  left=3cm,  right=2.5cm
]{geometry}

% --- Cores ---
\usepackage[dvipsnames]{xcolor}
\definecolor{primary}{HTML}{1A1A2E}
\definecolor{accent}{HTML}{E94560}
\definecolor{lightgray}{HTML}{F5F5F5}
\definecolor{codegray}{HTML}{2D2D2D}

% --- Links e PDF ---
\usepackage{hyperref}
\hypersetup{
  colorlinks = true,
  linkcolor  = accent,
  urlcolor   = accent,
  citecolor  = accent,
  pdftitle   = {Relatorio de Projeto: Podcastia},
  pdfauthor  = {Afonso Bitoque, Maria Neto, Duarte Cunha, José Tico}
}

% --- Gráficos e Figuras ---
\usepackage{graphicx}
\usepackage{float}
\usepackage{caption}
\captionsetup{font=small, labelfont=bf}

% --- Tabelas ---
\usepackage{booktabs}
\usepackage{longtable}
\usepackage{array}
\usepackage{tabularx}
\newcolumntype{L}[1]{>{\raggedright\arraybackslash}p{#1}}

% --- Código Fonte ---
\usepackage{listings}
\lstset{
  backgroundcolor = \color{lightgray},
  basicstyle      = \small\ttfamily,
  keywordstyle    = \color{accent}\bfseries,
  commentstyle    = \color{gray}\itshape,
  stringstyle     = \color{olive},
  numbers         = left,
  numberstyle     = \tiny\color{gray},
  numbersep       = 8pt,
  frame           = single,
  framerule       = 0.4pt,
  breaklines      = true,
  captionpos      = b,
  showstringspaces= false,
  tabsize         = 4
}
\lstdefinelanguage{bash}{
  keywords={curl,echo,cd,export,bash,git},
  sensitive=false,
  morecomment=[l]{\#},
  morestring=[b]"
}

% --- Cabeçalhos e Rodapés ---
\usepackage{fancyhdr}
\pagestyle{fancy}
\fancyhf{}
\fancyhead[L]{\small\textcolor{gray}{Podcastia: Relat\'{o}rio de Projeto}}
\fancyhead[R]{\small\textcolor{gray}{\leftmark}}
\fancyfoot[C]{\thepage}
\renewcommand{\headrulewidth}{0.4pt}

% --- Secções com cor ---
\usepackage{titlesec}
\titleformat{\chapter}[display]
  {\normalfont\huge\bfseries\color{primary}}
  {\chaptertitlename\ \thechapter}{20pt}{\Huge}
\titleformat{\section}
  {\normalfont\Large\bfseries\color{primary}}
  {\thesection}{1em}{}
\titleformat{\subsection}
  {\normalfont\large\bfseries\color{primary!80}}
  {\thesubsection}{1em}{}

% --- Caixas de destaque ---
\usepackage{tcolorbox}
\tcbuselibrary{skins}
\newtcolorbox{destaque}[1][]{
  enhanced,
  colback=accent!8,
  colframe=accent,
  boxrule=0.8pt,
  arc=4pt,
  fonttitle=\bfseries,
  #1
}
\newtcolorbox{infobox}[1][]{
  enhanced,
  colback=primary!5,
  colframe=primary,
  boxrule=0.8pt,
  arc=4pt,
  #1
}

% --- Misc ---
\usepackage{enumitem}
\usepackage{parskip}
\setlength{\parindent}{0pt}
\setlength{\parskip}{6pt}
\usepackage{emptypage}
\usepackage{pdfpages}

% ============================================================
\begin{document}

% ============================================================
%  CAPA
% ============================================================
\begin{titlepage}
  \pagecolor{primary}
  \color{white}
  \centering
  \vspace*{3cm}

  {\fontsize{48}{56}\selectfont\bfseries Podcastia}\\[0.4cm]
  {\large\color{accent} Plataforma de Podcasts Gerados por Inteligência Artificial}

  \vspace{2cm}
  \rule{10cm}{1pt}
  \vspace{1.5cm}

  {\large\bfseries Relatório de Projeto}\\[0.5cm]
  {\normalsize Licenciatura em Engenharia Informática}\\[0.2cm]
  {\normalsize Projeto de Desenvolvimento de Software}

  \vspace{2.5cm}

  {\large
    Afonso Bitoque\\[0.3cm]
    Maria Neto\\[0.3cm]
    Duarte Cunha\\[0.3cm]
    José Tico
  }

  \vfill
  {\normalsize Maio de 2026}
  \vspace{1cm}
\end{titlepage}
\nopagecolor

% ============================================================
%  PÁGINA DE ROSTO
% ============================================================
\newpage
\thispagestyle{empty}
\vspace*{2cm}

\begin{center}
  {\LARGE\bfseries\color{primary} Podcastia}\\[0.3cm]
  {\large Plataforma de Podcasts Gerados por Inteligência Artificial}\\[2cm]

  \begin{tabular}{ll}
    \toprule
    \textbf{Projeto}      & Podcastia \\
    \textbf{Versão}       & 1.0.0 \\
    \textbf{Data}         & Maio de 2026 \\
    \textbf{Repositório}  & \url{https://github.com/AfonsoBitoque/podcastia} \\
    \midrule
    \textbf{Autores}      & Afonso Bitoque \\
                          & Maria Neto \\
                          & Duarte Cunha \\
                          & José Tico \\
    \bottomrule
  \end{tabular}
\end{center}

\vfill

\begin{infobox}[title=Resumo]
O \textbf{Podcastia} é uma plataforma web de podcasts gerados automaticamente por
Inteligência Artificial. Qualquer utilizador registado pode criar um podcast sobre
qualquer tema: o sistema gera o guião com a API Google Gemini e sintetiza o áudio
em português europeu usando \texttt{edge-tts}. Para além da geração de conteúdo,
a plataforma oferece funcionalidades sociais completas, incluindo amizades, chat privado
em tempo real, playlists personalizadas, feed de recomendação e painel de
administração com analytics.
\end{infobox}

% ============================================================
%  ÍNDICE
% ============================================================
\newpage
\tableofcontents
\newpage

% ============================================================
\chapter{Introdução}
% ============================================================

\section{Contexto e Motivação}

O consumo de podcasts tem registado um crescimento exponencial nos últimos anos,
tornando-se um dos formatos de media mais populares a nível mundial. No entanto,
a criação de conteúdo de qualidade continua a exigir recursos significativos:
equipamento de gravação, edição de áudio e tempo de produção elevado.

O presente projeto explora a possibilidade de democratizar a criação de podcasts
recorrendo a Inteligência Artificial. O \textbf{Podcastia} permite que qualquer
utilizador, sem conhecimentos técnicos de produção áudio, crie um podcast completo
bastando indicar um tema. O sistema trata automaticamente da geração do guião e
da síntese de voz, entregando um ficheiro MP3 pronto a ouvir.

\section{Objetivos}

Os objetivos do projeto dividem-se em duas categorias:

\subsection{Objetivos Principais}
\begin{itemize}
  \item Desenvolver uma plataforma web funcional de criação e partilha de podcasts gerados por IA;
  \item Implementar um sistema de autenticação seguro e stateless baseado em JWT;
  \item Criar um motor de recomendação de conteúdo personalizado;
  \item Disponibilizar funcionalidades sociais (amizades, chat em tempo real, playlists);
  \item Fornecer um painel de administração com analytics e ferramentas de moderação.
\end{itemize}

\subsection{Objetivos Secundários}
\begin{itemize}
  \item Integrar consumo automático de feeds RSS de fontes externas;
  \item Implementar playlists diárias geradas automaticamente por um agendador;
  \item Assegurar qualidade de código através de testes automatizados e pipelines de CI/CD;
  \item Documentar a arquitetura utilizando o modelo C4 com diagramas em Mermaid.js.
\end{itemize}

\section{Estrutura do Relatório}

Este documento está organizado da seguinte forma: o Capítulo~\ref{chap:requisitos}
descreve os requisitos funcionais e não funcionais; o Capítulo~\ref{chap:arquitetura}
apresenta a arquitetura do sistema; o Capítulo~\ref{chap:tecnologias} lista as
tecnologias utilizadas; o Capítulo~\ref{chap:funcionalidades} detalha as
funcionalidades implementadas; os Capítulos~\ref{chap:seguranca}
e~\ref{chap:testes} abordam a segurança e os testes; o
Capítulo~\ref{chap:qualidade} analisa a qualidade e os problemas conhecidos;
o Capítulo~\ref{chap:gestao} apresenta a gestão do projeto com base nos dados do Jira;
o Capítulo~\ref{chap:instalacao} fornece instruções de instalação; e o
Capítulo~\ref{chap:conclusao} apresenta as conclusões e o trabalho futuro.

% ============================================================
\chapter{Requisitos do Sistema}
\label{chap:requisitos}
% ============================================================

\section{Requisitos Funcionais}

O sistema foi desenhado para satisfazer os seguintes requisitos funcionais, todos
implementados na versão atual:

\begin{longtable}{L{1cm} L{7cm} L{3cm}}
  \toprule
  \textbf{ID} & \textbf{Requisito} & \textbf{Estado} \\
  \midrule
  \endfirsthead
  \toprule
  \textbf{ID} & \textbf{Requisito} & \textbf{Estado} \\
  \midrule
  \endhead
  RF01 & Registo e autenticação de utilizadores (email ou username\#tag) & \textcolor{OliveGreen}{\checkmark\ Implementado} \\
  RF02 & Geração de podcasts por IA a partir de um tema & \textcolor{OliveGreen}{\checkmark\ Implementado} \\
  RF03 & Streaming e download de ficheiros de áudio MP3 & \textcolor{OliveGreen}{\checkmark\ Implementado} \\
  RF04 & Feed personalizado baseado em preferências do utilizador & \textcolor{OliveGreen}{\checkmark\ Implementado} \\
  RF05 & Sistema de amizades (pedido, aceitar, rejeitar, bloquear) & \textcolor{OliveGreen}{\checkmark\ Implementado} \\
  RF06 & Chat privado em tempo real entre utilizadores & \textcolor{OliveGreen}{\checkmark\ Implementado} \\
  RF07 & Criação e gestão de playlists personalizadas & \textcolor{OliveGreen}{\checkmark\ Implementado} \\
  RF08 & Playlist diária gerada automaticamente por agendador & \textcolor{OliveGreen}{\checkmark\ Implementado} \\
  RF09 & Sistema de favoritos para podcasts & \textcolor{OliveGreen}{\checkmark\ Implementado} \\
  RF10 & Pesquisa unificada de utilizadores e podcasts & \textcolor{OliveGreen}{\checkmark\ Implementado} \\
  RF11 & Perfil de utilizador com imagem e processo de onboarding & \textcolor{OliveGreen}{\checkmark\ Implementado} \\
  RF12 & Painel de administração com analytics e moderação & \textcolor{OliveGreen}{\checkmark\ Implementado} \\
  RF13 & Consumo automático de feeds RSS externos & \textcolor{OliveGreen}{\checkmark\ Implementado} \\
  RF14 & Exportação de dados em CSV e PDF pelo administrador & \textcolor{OliveGreen}{\checkmark\ Implementado} \\
  RF15 & Reações com emojis a mensagens de chat & \textcolor{OliveGreen}{\checkmark\ Implementado} \\
  RF16 & Download de playlist completa em formato ZIP & \textcolor{OliveGreen}{\checkmark\ Implementado} \\
  RF17 & Registo de progresso de escuta de podcasts & \textcolor{OliveGreen}{\checkmark\ Implementado} \\
  RF18 & Moderação de conteúdo (ocultar, destacar, marcar explícito) & \textcolor{OliveGreen}{\checkmark\ Implementado} \\
  \bottomrule
  \caption{Requisitos funcionais do sistema Podcastia}
\end{longtable}

\section{Requisitos Não Funcionais}

\begin{longtable}{L{1cm} L{5cm} L{5.5cm}}
  \toprule
  \textbf{ID} & \textbf{Requisito} & \textbf{Solução Aplicada} \\
  \midrule
  \endfirsthead
  \toprule
  \textbf{ID} & \textbf{Requisito} & \textbf{Solução Aplicada} \\
  \midrule
  \endhead
  RNF01 & Autenticação stateless e segura & JWT HS512, validade 24h, BCrypt \\
  RNF02 & Comunicação em tempo real & WebSocket STOMP com JWT no handshake \\
  RNF03 & Filtros de feed dinâmicos & JPA Specification (sem SQL concatenado) \\
  RNF04 & Proteção contra abuso no chat & Rate limiting 20\,msg/min e 100\,msg/h \\
  RNF05 & Imagens de perfil controladas & Validação $\leq$\,5\,MB, redimensionamento 400$\times$400\,px \\
  RNF06 & Cache de recomendações & \texttt{ConcurrentHashMap} in-memory (24\,h) \\
  RNF07 & Qualidade de código & Checkstyle (Google style) + OWASP \\
  RNF08 & Portabilidade de base de dados & JPA/Hibernate (H2 dev $\rightarrow$ PostgreSQL prod) \\
  \bottomrule
  \caption{Requisitos não funcionais do sistema Podcastia}
\end{longtable}

% ============================================================
\chapter{Arquitetura do Sistema}
\label{chap:arquitetura}
% ============================================================

\section{Visão Geral}

A arquitetura do Podcastia segue o modelo \textbf{C4} (Context, Containers,
Components, Code), documentado em quatro níveis com diagramas Mermaid.js.
O sistema assenta numa arquitectura cliente-servidor clássica, com clara separação
entre o frontend \emph{Single-Page Application} e o backend REST API, complementada
por comunicação em tempo real via WebSocket STOMP.

\section{Contexto do Sistema (C4 Nível 1)}

O Podcastia interage com quatro sistemas externos:

\begin{itemize}
  \item \textbf{Google Gemini API}: geração do guião narrativo do podcast em português;
  \item \textbf{edge-tts (Python)}: síntese de texto em áudio MP3 com a voz \texttt{pt-PT-RaquelNeural};
  \item \textbf{Fontes RSS Externas}: feeds de parceiros (Observador, Público Desporto, TechCrunch, BBC News) consumidos a cada duas horas;
  \item \textbf{Serviço de Email}: envio de notificações e relatórios (actualmente em modo \emph{stub} para desenvolvimento).
\end{itemize}

\section{Containers (C4 Nível 2)}

O sistema é composto por quatro containers principais, descritos na
Tabela~\ref{tab:containers}.

\begin{table}[H]
  \centering
  \begin{tabularx}{\textwidth}{L{3cm} L{3.5cm} X}
    \toprule
    \textbf{Container} & \textbf{Tecnologia} & \textbf{Responsabilidade} \\
    \midrule
    Frontend SPA       & React 19 + Vite 7   & Interface web; comunica com o backend via REST e WebSocket STOMP \\
    Backend API        & Spring Boot 3.4 + Java 17 & Lógica de negócio, autenticação JWT, agendadores e integrações \\
    Base de Dados      & H2 / SQL            & Persistência de utilizadores, podcasts, playlists, mensagens e logs \\
    Sistema de Ficheiros & Disco local        & Armazenamento de ficheiros MP3 e imagens de perfil \\
    \bottomrule
  \end{tabularx}
  \caption{Containers do sistema Podcastia}
  \label{tab:containers}
\end{table}

\section{Componentes do Backend (C4 Nível 3)}

O backend está organizado em quatro camadas bem definidas, seguindo os princípios
da arquitectura em camadas e da inversão de dependências.

\subsection{Camada de Segurança}

A segurança do sistema é implementada de forma transversal através de quatro
componentes:

\begin{description}
  \item[\texttt{SecurityConfig}:] Configura o Spring Security, define os endpoints públicos, as regras de CORS e integra o filtro JWT.
  \item[\texttt{JwtAuthenticationFilter}:] Intercepta todos os pedidos HTTP, valida o token Bearer e popula o \texttt{SecurityContext}.
  \item[\texttt{JwtUtil}:] Gera e valida tokens JWT HS512 com validade de 24 horas, transportando os claims \texttt{sub} (email), \texttt{id} e \texttt{type}.
  \item[\texttt{JwtWebSocketHandshakeInterceptor}:] Autentica as ligações WebSocket verificando o token JWT no handshake HTTP inicial.
\end{description}

\subsection{Camada de Controllers (17 controllers)}

Os controllers são responsáveis por receber os pedidos HTTP/WebSocket, validar
os dados de entrada e delegar o processamento para a camada de serviços. Alguns
exemplos relevantes:

\begin{itemize}
  \item \texttt{PodcastGenerationController} (\texttt{/api/podcasts}): geração de podcasts por IA, streaming e download de áudio;
  \item \texttt{FeedController} (\texttt{/api/home}): feed filtrado com múltiplos parâmetros combinados;
  \item \texttt{ChatWebSocketController} (STOMP \texttt{/app/chat.*}): mensagens em tempo real;
  \item \texttt{AdminController} (\texttt{/api/admin}): painel de administração, exclusivo a utilizadores com papel \texttt{USER\_ADMIN}.
\end{itemize}

\subsection{Camada de Serviços (13 serviços)}

Os serviços encapsulam toda a lógica de negócio. Destacam-se:

\begin{description}
  \item[\texttt{PodcastGenerationService}:] Implementa o pipeline completo de geração, desde a chamada à Gemini API, passando pela limpeza do guião e síntese TTS, até à persistência do MP3.
  \item[\texttt{RecommendationService}:] Calcula um feed personalizado com base nos pontos de cada categoria (DESPORTO, POLÍTICA, FINANÇAS, GERAL), mantendo uma cache de 24 horas por utilizador.
  \item[\texttt{ChatMessageServiceImpl}:] Gere o ciclo completo de mensagens, incluindo envio, entrega, leitura, reações e eliminação, com rate limiting e blacklist de links.
  \item[\texttt{UserRelationshipServiceImpl}:] Implementa o ciclo de amizades com cooldown de 7 dias após rejeição e gestão de bloqueios.
  \item[\texttt{DailyPlaylistService}:] Gera automaticamente uma playlist diária personalizada para cada utilizador com base nas suas preferências.
\end{description}

\subsection{Camada de Repositórios (13 repositórios)}

Todos os repositórios estendem \texttt{JpaRepository} do Spring Data JPA.
Para operações complexas utilizam-se queries JPQL anotadas com \texttt{@Query},
evitando SQL manual. O \texttt{PodcastRepository} implementa ainda
\texttt{JpaSpecificationExecutor} para suportar os filtros dinâmicos do feed.

\section{Pipeline de Geração de Podcasts}

O fluxo de geração de um podcast, desde o pedido do utilizador até à entrega
do áudio, é ilustrado de seguida:

\begin{enumerate}
  \item O utilizador envia \texttt{POST /api/podcasts/generate} com o tema e categorias desejadas;
  \item O \texttt{PodcastGenerationService} constrói um \emph{prompt} contextualizado e chama a Google Gemini API;
  \item A API devolve um guião narrativo em português europeu;
  \item O guião é limpo (remoção de formatação Markdown) e passado ao processo Python \texttt{edge-tts} com a voz \texttt{pt-PT-RaquelNeural};
  \item O ficheiro MP3 resultante é guardado em \texttt{generated-podcasts/};
  \item Os metadados do podcast são persistidos na base de dados;
  \item A resposta inclui o \texttt{podcastId} e a URL de \emph{streaming} imediatamente utilizável.
\end{enumerate}

\begin{destaque}[title=Nota sobre desempenho]
  O processo de geração é síncrono e pode demorar entre 30 e 60 segundos, dependendo
  da dimensão do guião gerado e da velocidade da rede. Para podcasts com 3 a 5 minutos
  de duração, o guião gerado contém tipicamente entre 450 e 750 palavras.
\end{destaque}

% ============================================================
\chapter{Tecnologias Utilizadas}
\label{chap:tecnologias}
% ============================================================

\section{Backend}

\begin{table}[H]
  \centering
  \begin{tabularx}{\textwidth}{L{4cm} L{2cm} X}
    \toprule
    \textbf{Tecnologia} & \textbf{Versão} & \textbf{Utilização} \\
    \midrule
    Java               & 17        & Linguagem principal do backend \\
    Spring Boot        & 3.4.3     & Framework principal (web, security, data, websocket) \\
    Spring Security    & 6.x       & Autenticação JWT e autorização por papéis \\
    Spring Data JPA    & 3.x       & Persistência, queries JPQL e Specification \\
    Hibernate          & 6.x       & ORM e DDL automático \\
    H2 Database        & n/a       & Base de dados para desenvolvimento e testes \\
    JJWT               & 0.11.5    & Geração e validação de tokens JWT HS512 \\
    ROME Library       & 2.1.0     & Parsing de feeds RSS/Atom \\
    SpringDoc OpenAPI  & 2.8.5     & Documentação automática da API (Swagger UI) \\
    Python edge-tts    & 7.2.8     & Síntese de texto em áudio MP3 com voz PT \\
    Google Gemini API  & n/a       & Geração de guiões narrativos com LLM \\
    Maven              & 3.x       & Gestão de dependências e automação de build \\
    \bottomrule
  \end{tabularx}
  \caption{Tecnologias do backend}
\end{table}

\section{Frontend}

\begin{table}[H]
  \centering
  \begin{tabularx}{\textwidth}{L{4cm} L{2cm} X}
    \toprule
    \textbf{Tecnologia} & \textbf{Versão} & \textbf{Utilização} \\
    \midrule
    React              & 19.2      & Framework de UI (componentes, hooks, routing) \\
    Vite               & 7.3       & Bundler e servidor de desenvolvimento HMR \\
    React Router       & 7.13      & Routing client-side (SPA) \\
    react-hot-toast    & 2.4       & Notificações visuais não-intrusivas \\
    WebSocket STOMP    & n/a       & Chat privado em tempo real \\
    Service Worker     & n/a       & Suporte offline (PWA) \\
    Vitest             & 4.x       & Testes unitários de componentes \\
    Playwright         & 1.45      & Testes End-to-End no browser \\
    ESLint             & 9.x       & Análise estática de código JavaScript \\
    Prettier           & 3.3       & Formatação automática de código \\
    \bottomrule
  \end{tabularx}
  \caption{Tecnologias do frontend}
\end{table}

% ============================================================
\chapter{Funcionalidades Implementadas}
\label{chap:funcionalidades}
% ============================================================

\section{Autenticação e Registo}

O sistema suporta dois modos de autenticação, ambos devolvendo um token JWT
com validade de 24 horas:

\begin{itemize}
  \item \textbf{Por email}: \texttt{\{"identifier": "user@email.com", "password": "..."\}}
  \item \textbf{Por username\#tag}: \texttt{\{"identifier": "nome\#0000", "password": "..."\}}
\end{itemize}

O sistema de tags (números de 0000 a 9999) permite que múltiplos utilizadores
partilhem o mesmo nome de utilizador sem ambiguidade, de forma análoga ao modelo
utilizado pelo Discord. A tag é atribuída automaticamente no momento do registo.

O \textbf{onboarding} pós-registo recolhe os tópicos de interesse do utilizador,
que alimentam o motor de recomendação desde a primeira sessão.

\section{Geração de Podcasts por Inteligência Artificial}

Esta é a funcionalidade central da plataforma. Um utilizador autenticado pode
gerar um podcast preenchendo apenas dois campos:

\begin{itemize}
  \item \textbf{Tema}: assunto do podcast em linguagem natural (por exemplo, \emph{"Curiosidades sobre o espaço sideral"});
  \item \textbf{Categorias}: uma ou mais de \texttt{GERAL}, \texttt{DESPORTO}, \texttt{POLITICA}, \texttt{FINANCAS}.
\end{itemize}

O processo é completamente automático: o sistema chama a Google Gemini API com
um \emph{prompt} estruturado, recebe um guião narrativo, passa-o ao
\texttt{edge-tts} com a voz \texttt{pt-PT-RaquelNeural} e devolve o MP3 pronto
a reproduzir.

\section{Feed Personalizado e Recomendação}

O feed principal suporta filtros combinados em tempo real:

\begin{itemize}
  \item \texttt{type=recommended}: podcasts ordenados por pontuação de relevância nas categorias do utilizador;
  \item \texttt{category=DESPORTO|POLITICA|FINANCAS|GERAL}: filtro por categoria;
  \item \texttt{max\_duration=30}: duração máxima em minutos;
  \item \texttt{hide\_played=true}: oculta podcasts já ouvidos;
  \item \texttt{shorts=true}: apenas podcasts curtos ($\leq$ 5 minutos);
  \item \texttt{is\_favorite=true}: apenas favoritos do utilizador.
\end{itemize}

O motor de recomendação mantém uma pontuação por categoria para cada utilizador.
Cada podcast ouvido incrementa os pontos da sua categoria, influenciando as
recomendações futuras de forma progressiva.

\section{Sistema Social}

\subsection{Amizades}

O sistema de amizades segue um modelo direcional com os seguintes estados:
\texttt{PEDIDO\_PENDENTE}, \texttt{AMIGO}, \texttt{BLOQUEADO},
\texttt{PEDIDO\_REJEITADO} e \texttt{CANCELADO}. Uma amizade aceite cria dois
registos na base de dados (um por sentido), garantindo a bidirecionalidade das
consultas. Após a rejeição de um pedido, existe um período de \emph{cooldown}
de 7 dias antes de ser possível enviar um novo pedido.

\subsection{Chat Privado em Tempo Real}

O chat utiliza o protocolo WebSocket STOMP, com autenticação JWT verificada no
\emph{handshake} HTTP inicial. As mensagens suportam três estados de entrega:
\texttt{SENT}, \texttt{DELIVERED} e \texttt{READ}. O histórico de conversas
é disponibilizado com paginação cursor-based (por \texttt{Instant}), evitando
os problemas de estabilidade do offset tradicional em feeds com inserções frequentes.

O sistema inclui mecanismos de proteção contra abuso:
\begin{itemize}
  \item Rate limiting: 20 mensagens por minuto e 100 por hora por utilizador;
  \item Blacklist de links: utilizadores que enviam links excessivamente sofrem bloqueio progressivo;
  \item Reações com emojis às mensagens.
\end{itemize}

\section{Playlists}

\subsection{Playlists Personalizadas}

Os utilizadores podem criar playlists com título, descrição e visibilidade
pública ou privada. É possível adicionar, remover e reordenar episódios, partilhar
playlists públicas com amigos e descarregar uma playlist completa como ficheiro
ZIP contendo todos os MP3.

\subsection{Playlist Diária Automática}

O agendador \texttt{DailyPlaylistScheduler} gera automaticamente, à meia-noite
(fuso horário \texttt{Europe/Lisbon}), uma playlist personalizada para cada
utilizador activo. O algoritmo selecciona até 20 podcasts com duração mínima
total de 30 minutos, ordenados por pontuação de relevância calculada com base
nos pontos de categoria do utilizador.

\section{Painel de Administração}

O painel de administração, acessível exclusivamente a utilizadores com o papel
\texttt{USER\_ADMIN}, disponibiliza:

\begin{itemize}
  \item \textbf{Analytics em tempo real}: DAU, MAU, novos registos, top podcasts, uso semanal e mensal, saúde do sistema;
  \item \textbf{Gestão de conteúdo}: ocultar, destacar, marcar como explícito e eliminar podcasts (com dupla confirmação);
  \item \textbf{Gestão de utilizadores}: reset de password, suspensão e eliminação de contas;
  \item \textbf{Auditoria}: log histórico de todas as acções administrativas com timestamp e resultado;
  \item \textbf{Exportação}: relatórios em CSV e PDF, com geração assíncrona por \emph{job}.
\end{itemize}

% ============================================================
\chapter{Segurança}
\label{chap:seguranca}
% ============================================================

\section{Modelo de Segurança}

A segurança do Podcastia é implementada em múltiplas camadas, seguindo uma
abordagem de defesa em profundidade (\emph{defense in depth}):

\subsection{Autenticação JWT}

O sistema utiliza tokens JWT com o algoritmo HMAC-SHA512 (HS512), com validade
de 24 horas. Cada token transporta os seguintes \emph{claims}: \texttt{sub}
(email do utilizador), \texttt{id} (identificador único) e \texttt{type}
(papel do utilizador). A transmissão é feita exclusivamente no header
\texttt{Authorization: Bearer <token>}, sem armazenamento em cookies (eliminando
o risco de CSRF).

\subsection{Proteção de Passwords}

As passwords são armazenadas exclusivamente como hashes BCrypt, com o fator de
custo padrão do Spring Security. Nenhum componente do sistema tem acesso à
password em texto simples após o registo.

\subsection{Controlo de Acesso}

O controlo de acesso é implementado em dois níveis:
\begin{itemize}
  \item \textbf{Declarativo}: os endpoints de administração usam a anotação \texttt{@PreAuthorize("hasRole('USER\_ADMIN')")} do Spring Security, verificada pelo AOP antes da execução da lógica de negócio;
  \item \textbf{Programático}: operações sobre recursos (por exemplo, editar ou eliminar uma playlist) verificam programaticamente se o utilizador autenticado é o proprietário do recurso.
\end{itemize}

\section{Vulnerabilidades Conhecidas}

\begin{table}[H]
  \centering
  \begin{tabularx}{\textwidth}{L{5cm} L{2cm} X}
    \toprule
    \textbf{Vulnerabilidade} & \textbf{Severidade} & \textbf{Estado e Mitigação} \\
    \midrule
    Chave JWT \emph{hardcoded} em \texttt{JwtUtil} & Alta & Identificada; deve ser movida para variável de ambiente antes do deploy em produção \\
    Logging de tokens JWT no frontend & Alta & Identificada; os \texttt{console.log} devem ser removidos em produção \\
    Mensagens de erro com detalhes internos & Média & Identificada; as mensagens devem ser sanitizadas \\
    CORS configurado para \texttt{localhost} & Baixa & Resolvida para produção; restrito a origens específicas no \texttt{SecurityConfig} \\
    \bottomrule
  \end{tabularx}
  \caption{Vulnerabilidades de segurança identificadas}
\end{table}

% ============================================================
\chapter{Testes}
\label{chap:testes}
% ============================================================

\section{Estratégia de Testes}

O projeto segue uma pirâmide de testes com cobertura nas camadas de serviço e
integração. A suite inclui 12 ficheiros de teste no backend, complementados por
testes de componente (Vitest) e testes E2E (Playwright) no frontend.

\section{Suite de Testes Backend}

\begin{longtable}{L{5.5cm} L{2.5cm} L{5.5cm}}
  \toprule
  \textbf{Ficheiro} & \textbf{Tipo} & \textbf{O Que Cobre} \\
  \midrule
  \endfirsthead
  \toprule
  \textbf{Ficheiro} & \textbf{Tipo} & \textbf{O Que Cobre} \\
  \midrule
  \endhead
  \texttt{ServidorApplicationTests}       & Context Load  & Arranque correto do contexto Spring Boot \\
  \texttt{UserRegistrationTest}           & Integração    & Criação de utilizador, conflitos, validações \\
  \texttt{ChangePasswordIntegrationTest}  & Integração    & Alteração de password com verificação BCrypt \\
  \texttt{AuthIntegrationTest}            & Integração    & Login email/username\#tag, JWT, endpoints protegidos \\
  \texttt{UserRelationshipServiceTest}    & Unitário      & Pedidos de amizade, cooldown, bloqueios (Mockito) \\
  \texttt{UserRelationIntegrationTest}    & Integração    & Ciclo completo de amizades via API REST \\
  \texttt{FeedIntegrationTest}            & Integração    & Filtros do feed (categoria, shorts, favoritos) \\
  \texttt{PlaylistIntegrationTest}        & Integração    & CRUD de playlists, episódios, permissões \\
  \texttt{ChatMessageServiceTest}         & Unitário      & Envio, rate limiting, links, permissões \\
  \texttt{ChatReactionServiceTest}        & Unitário      & Adicionar, remover e actualizar reações \\
  \texttt{RssServiceTest}                 & Integração    & Consumo RSS, deduplicação, fallbacks \\
  \texttt{OpenApiIntegrationTest}         & Integração    & Swagger UI e \texttt{/v3/api-docs} sem autenticação \\
  \bottomrule
  \caption{Suite de testes automatizados do backend}
\end{longtable}

\section{Execução dos Testes}

\begin{lstlisting}[language=bash, caption={Comandos para executar os testes}]
# Todos os testes
cd servidor && ./mvnw test

# Teste especifico
./mvnw test -Dtest=UserRegistrationTest

# Validacao completa (build + testes + checkstyle + OWASP)
./mvnw verify
\end{lstlisting}

% ============================================================
\chapter{Pipeline de Integração Contínua}
% ============================================================

\section{GitHub Actions}

O projeto utiliza duas pipelines de CI configuradas como GitHub Actions:

\subsection{Backend CI (\texttt{maven-ci.yml})}

Activa em qualquer \emph{push} ou em \emph{pull requests} para \texttt{main}/\texttt{master}.
O fluxo provisiona Ubuntu, instala o JDK 17 (Temurin) com cache Maven e executa
\texttt{mvn -B verify}, que compila o projecto, executa todos os testes,
verifica o estilo de código com o Checkstyle e audita as dependências com o
OWASP Dependency Check.

\subsection{Frontend CI (\texttt{frontend-ci.yml})}

Activa em alterações à pasta \texttt{frontend/}. O fluxo instala o Node.js 20
com cache npm, executa \texttt{npm run lint} (ESLint), \texttt{npm run test:ci}
(Vitest) e os testes E2E com Playwright. O relatório Playwright é publicado como
artefacto da pipeline com retenção de 30 dias.

% ============================================================
\chapter{Análise de Qualidade e Problemas Conhecidos}
\label{chap:qualidade}
% ============================================================

\section{Métricas do Código}

\begin{table}[H]
  \centering
  \begin{tabular}{lr}
    \toprule
    \textbf{Métrica} & \textbf{Valor} \\
    \midrule
    Classes Java (backend)   & $\approx$ 91 \\
    Controllers              & 17 \\
    Serviços                 & 13 \\
    Entidades / Modelos      & 16 \\
    Repositórios JPA         & 13 \\
    Ficheiros de teste       & 12 \\
    Ficheiros JSX (frontend) & $\approx$ 35 \\
    Hooks customizados React & 3 \\
    \bottomrule
  \end{tabular}
  \caption{Métricas de dimensão do projecto}
\end{table}

\section{Bugs Resolvidos}

Durante o desenvolvimento foram identificados e resolvidos vários bugs:

\begin{itemize}
  \item \textbf{Loop de redirecionamento pós-onboarding} (crítico): resolvido com evento \texttt{auth-change} para forçar a re-renderização reactiva do \texttt{ProtectedRoute};
  \item \textbf{Endpoint \texttt{GET /api/podcasts} inexistente} (crítico): adicionado o método \texttt{getAllPublicPodcasts()} com filtro \texttt{hidden=false};
  \item \textbf{Recursos estáticos com 403} (crítico): criado \texttt{StaticResourceConfig} com handlers para \texttt{/images/**} e \texttt{/audio/**};
  \item \textbf{\texttt{UnsupportedOperationException} no onboarding} (crítico): \texttt{stream().toList()} retorna lista imutável; corrigido com conversão para \texttt{ArrayList};
  \item \textbf{\texttt{window.confirm()} bloqueado no browser} (médio): o diálogo nativo é bloqueado silenciosamente por políticas do browser moderno, tendo sido substituído por modal React interno.
\end{itemize}

\section{Dívida Técnica}

\begin{table}[H]
  \centering
  \begin{tabularx}{\textwidth}{X L{2.5cm}}
    \toprule
    \textbf{Issue} & \textbf{Prioridade} \\
    \midrule
    Chave JWT \emph{hardcoded} em \texttt{JwtUtil} & Alta \\
    \texttt{removeFavorite} usa \texttt{findAll()} (ineficiente) & Média \\
    Estado de autenticação duplicado no \texttt{localStorage} & Média \\
    Sem \emph{lazy loading} de componentes React & Baixa \\
    \texttt{AuthUserController} mistura autenticação e onboarding (SRP) & Baixa \\
    \bottomrule
  \end{tabularx}
  \caption{Dívida técnica identificada}
\end{table}

% ============================================================
\chapter{Gestão do Projeto}
\label{chap:gestao}
% ============================================================

\section{Metodologia}

O desenvolvimento do Podcastia seguiu uma metodologia ágil baseada em Scrum,
com sprints semanais geridas através da plataforma \textbf{Jira} (projeto
\texttt{POD} em \url{https://ualg-team-podcast.atlassian.net}). O backlog foi
constituído por histórias de utilizador (User Stories) e tarefas técnicas,
estimadas em story points. No total, o projeto registou \textbf{111 pendências}
distribuídas por 12 sprints, todas concluídas.

\section{Equipa e Distribuição de Trabalho}

A equipa era composta por quatro elementos, cada um com áreas de responsabilidade
primárias, embora com colaboração transversal entre todos:

\begin{table}[H]
  \centering
  \begin{tabularx}{\textwidth}{L{4cm} X}
    \toprule
    \textbf{Elemento} & \textbf{Principais Responsabilidades} \\
    \midrule
    Afonso Bitoque       & Backend (motor de recomendações, algoritmo de feed, pipeline de geração de podcasts, onboarding, endpoint da homepage, cache) \\
    Maria Neto           & Frontend (UI/UX, player, pesquisa, perfil, sistema de mensagens, playlists, páginas de autenticação, polimento visual) \\
    Duarte Cunha         & Backend (autenticação JWT, sistema de amizades, testes de integração, pipeline CI/CD, consumo RSS, Swagger/OpenAPI) \\
    José Tico            & Backend/Frontend (playlists, sistema de mensagens, reações, pesquisa com filtros, funcionalidades de perfil, hot fixes) \\
    \bottomrule
  \end{tabularx}
  \caption{Distribuição de responsabilidades por elemento da equipa}
\end{table}

\section{Progresso por Sprint}

A tabela seguinte resume o progresso do projeto sprint a sprint, com base nos
dados exportados do Jira em 24 de maio de 2026.

\begin{table}[H]
  \centering
  \begin{tabular}{L{2.5cm} r r r L{6cm}}
    \toprule
    \textbf{Sprint} & \textbf{Issues} & \textbf{Concluídos} & \textbf{Story Points} & \textbf{Foco Principal} \\
    \midrule
    Sprint 1  & 7  & 7  & 21,0 & Configuração do projeto Spring Boot, modelos de dados e endpoints base \\
    Sprint 2  & 4  & 4  & 3,0  & Backend de autenticação: criação de conta, login, gestão de contas \\
    Sprint 3  & 4  & 4  & 3,5  & Frontend de autenticação, testes unitários, pipeline de CI/CD \\
    Sprint 4  & 4  & 4  & 5,5  & Migração para JWT, logout, cabeçalho com nome de utilizador \\
    Sprint 5  & 11 & 11 & 11,2 & Recomendações, cache, playlist, foto de perfil, OpenAPI, retoma de reprodução \\
    Sprint 6  & n/a & n/a & n/a  & Issues partilhadas com outros sprints (ver Sprint 5/7) \\
    Sprint 7  & 9  & 9  & 13,0 & Sistema de amizades (backend), testes, sidebar de detalhe do podcast \\
    Sprint 8  & 5  & 5  & 5,0  & Endpoints de amizade (aceitar, recusar, listar, remover) \\
    Sprint 9  & 3  & 3  & 3,5  & Sistema de mensagens, reações, playlist diária automática \\
    Sprint 10 & 8  & 8  & 10,0 & Dashboard admin, gestão de conteúdo, frontend de mensagens e reações \\
    Sprint 11 & 8  & 8  & 3,5  & Hot fixes, melhorias de UX, notificações, pesquisa sem login \\
    Sprint 12 & 9  & 9  & 9,0  & Polimento visual, testes frontend, correção de bugs de áudio e pesquisa \\
    \midrule
    \textbf{Total} & \textbf{111} & \textbf{111} & \textbf{88,2} & \\
    \bottomrule
  \end{tabular}
  \caption{Resumo do progresso por sprint (fonte: Jira, 24/05/2026)}
\end{table}

\section{Destaques por Sprint}

\subsection{Sprints 1 a 4: Fundações}

Os primeiros quatro sprints estabeleceram as fundações do sistema: a estrutura
do projeto Spring Boot, o modelo de dados, os endpoints de autenticação (registo,
login, alteração de password) e a migração para autenticação JWT. Em paralelo,
foi criada a interface de autenticação no frontend e configurada a pipeline de
CI/CD com GitHub Actions.

\subsection{Sprint 5: Motor Central}

O quinto sprint foi o mais denso em termos técnicos, com 11 tarefas concluídas.
Foram implementados o algoritmo de recomendações por histórico, a cache de feed,
o endpoint agregador da homepage, a retoma de reprodução e a criação de playlists.

\subsection{Sprints 7 e 8: Sistema Social}

Os sprints 7 e 8 focaram-se no sistema de amizades: implementação dos endpoints
para enviar, aceitar, recusar, cancelar e remover ligações de amizade, bem como
os respectivos testes unitários e de integração. No frontend, foi desenvolvida
a sidebar de detalhes do podcast e a navegação directa para o ecrã de player.

\subsection{Sprints 9 e 10: Comunicação e Administração}

O sprint 9 introduziu o sistema de mensagens privadas e de reações, bem como a
automatização da playlist diária. O sprint 10 completou o painel de administração
(dashboard de estatísticas e gestão de conteúdo) e o frontend das funcionalidades
de comunicação.

\subsection{Sprints 11 e 12: Qualidade e Polimento}

Os dois últimos sprints foram dedicados a correcções de bugs, melhorias de
experiência de utilizador e polimento visual. Destacam-se a correcção do estado
do áudio ao navegar entre rotas (Sprint 12), a implementação de testes de
frontend com Playwright e a refactorização das páginas de playlists, tendências
e homepage.

% ============================================================
\chapter{Guia de Instalação e Execução}
\label{chap:instalacao}
% ============================================================

\section{Pré-requisitos}

\begin{table}[H]
  \centering
  \begin{tabular}{lll}
    \toprule
    \textbf{Requisito} & \textbf{Versão} & \textbf{Observação} \\
    \midrule
    Java JDK  & 17+  & Necessário para o backend \\
    Maven     & 3.8+ & Ou usar o wrapper \texttt{./mvnw} incluído \\
    Node.js   & 20+  & Necessário para o frontend \\
    Python    & 3.x  & Necessário para a síntese TTS \\
    pip       & n/a  & Para instalar o \texttt{edge-tts} \\
    \bottomrule
  \end{tabular}
  \caption{Pré-requisitos de instalação}
\end{table}

\section{Instalação}

\begin{lstlisting}[language=bash, caption={Passos de instalação}]
# 1. Clonar o repositorio
git clone https://github.com/AfonsoBitoque/podcastia.git
cd podcastia

# 2. Instalar dependencia Python para TTS
python3 -m pip install edge-tts --break-system-packages

# 3. Instalar dependencias do frontend
cd frontend && npm install && cd ..

# 4. Configurar a chave Gemini (necessario para geracao de podcasts)
echo "GEMINI_API_KEY=SUA_CHAVE_AQUI" > servidor/env.properties
\end{lstlisting}

\section{Execução}

\begin{lstlisting}[language=bash, caption={Iniciar o backend e o frontend}]
# Backend (porta 8080)
cd servidor && ./mvnw spring-boot:run

# Frontend (porta 5173) -- noutra janela de terminal
cd frontend && npm run dev
\end{lstlisting}

Após o arranque, o sistema estará disponível em:
\begin{itemize}
  \item \textbf{Frontend}: \url{http://localhost:5173}
  \item \textbf{API}: \url{http://localhost:8080}
  \item \textbf{Swagger UI}: \url{http://localhost:8080/swagger-ui.html}
  \item \textbf{H2 Console}: \url{http://localhost:8080/h2-console}
\end{itemize}

\begin{infobox}[title=Credenciais Padrão (DataSeeder)]
  O sistema cria automaticamente um utilizador administrador no primeiro arranque.\\[0.3cm]
  \textbf{Email:} \texttt{admin@podcastia.com} \quad
  \textbf{Password:} \texttt{admin} \quad
  \textbf{Papel:} \texttt{USER\_ADMIN}
\end{infobox}

% ============================================================
\chapter{Decisões Arquiteturais}
% ============================================================

As principais decisões de arquitetura do projecto estão resumidas na tabela
seguinte, com a justificação de cada escolha:

\begin{longtable}{L{3.5cm} L{3cm} L{7cm}}
  \toprule
  \textbf{Decisão} & \textbf{Escolha} & \textbf{Justificação} \\
  \midrule
  \endfirsthead
  \toprule
  \textbf{Decisão} & \textbf{Escolha} & \textbf{Justificação} \\
  \midrule
  \endhead
  Autenticação         & JWT stateless HS512 (24h)       & Sem estado no servidor, escalável horizontalmente \\
  Chat em tempo real   & WebSocket STOMP                 & Bidirecional, integra com Spring Messaging, suporte a subscrições por tópico \\
  ORM                  & JPA/Hibernate + H2              & DDL automático em desenvolvimento, migração para PostgreSQL/MySQL com uma linha de configuração \\
  Filtros de feed      & JPA Specification               & Queries dinâmicas seguras sem concatenação SQL \\
  Geração de áudio     & edge-tts via ProcessBuilder     & TTS gratuito com voz portuguesa nativa de alta qualidade \\
  Geração de guião     & Google Gemini API               & LLM com qualidade comprovada em português europeu \\
  Cache de feed        & ConcurrentHashMap (24h)         & Suficiente para instância única, substituível por Redis em cluster \\
  Playlists diárias    & @Scheduled cron à meia-noite    & Leve e não interactivo, sem necessidade de job queue externo \\
  Consumo RSS          & ROME Library + @Scheduled 2h    & Abstrai diferenças RSS\,1.0/2.0/Atom, biblioteca madura \\
  Segurança admin      & @PreAuthorize AOP               & Declarativo, verificado antes de chegar à lógica de negócio \\
  Paginação do chat    & Cursor-based (Instant)          & Estável em feeds com inserções frequentes, evita problemas de offset \\
  Tags de utilizador   & username\#0000-9999             & Permite nomes duplicados sem ambiguidade \\
  \bottomrule
  \caption{Decisões arquiteturais do projecto Podcastia}
\end{longtable}

% ============================================================
\chapter{Conclusão}
\label{chap:conclusao}
% ============================================================

\section{Síntese dos Resultados}

O projecto \textbf{Podcastia} foi desenvolvido e entregue como uma plataforma
web completa e funcional. Todos os 18 requisitos funcionais identificados foram
implementados, incluindo a funcionalidade central de geração de podcasts por
Inteligência Artificial, o sistema social (amizades, chat, playlists), o motor
de recomendação personalizado e o painel de administração.

A arquitetura adoptada demonstrou ser sólida e extensível: a separação em
camadas (segurança, controllers, serviços, repositórios) facilita a manutenção e
a adição de novas funcionalidades sem impacto nas camadas existentes. A escolha
do Spring Boot como framework principal, combinada com o Spring Security e o
Spring Data JPA, permitiu desenvolver um sistema robusto em tempo lectivo.

A qualidade do código foi assegurada através de pipelines de CI automáticas
(GitHub Actions), uma suite de 12 testes automatizados que cobre os fluxos
críticos, Checkstyle com as regras Google e auditoria de segurança OWASP das
dependências.

\section{Trabalho Futuro}

A versão actual do Podcastia constitui uma base sólida sobre a qual é possível
construir funcionalidades adicionais. As melhorias mais relevantes para trabalho
futuro incluem:

\begin{description}
  \item[Curto prazo:]
    Mover a chave JWT para variável de ambiente; implementar os serviços \texttt{NotificationService} e \texttt{EmailService} reais; adicionar endpoint de health check para as dependências externas.
  \item[Médio prazo:]
    Migrar a base de dados para PostgreSQL para produção; implementar \emph{lazy loading} de páginas React; substituir a cache em memória por Redis para arquitectura distribuída.
  \item[Longo prazo:]
    Suporte a múltiplas vozes TTS e idiomas; sistema de comentários em podcasts; integração com plataformas externas como Spotify e Apple Podcasts; deploy em cloud com Docker e CI/CD automático.
\end{description}

\section{Considerações Finais}

O Podcastia demonstra que é possível criar uma plataforma de media completa,
com funcionalidades sociais e conteúdo gerado por inteligência artificial, utilizando
tecnologias de código aberto ou gratuitas (Spring Boot, React, edge-tts) e um investimento
mínimo em serviços externos (Google Gemini API). O projecto alcançou todos os
objectivos propostos e constitui uma base técnica sólida para desenvolvimento futuro.

\vspace{2cm}
\begin{center}
  \rule{8cm}{0.4pt}\\[0.3cm]
  {\small Afonso Bitoque \quad|\quad Maria Neto \quad|\quad Duarte Cunha \quad|\quad José Tico}\\
  {\small Maio de 2026}
\end{center}

\end{document}
